% ============================================================================
%                         PAGE LAYOUT & GEOMETRY
% ============================================================================
% Geometry and layout per Udinus guide (BAB IV 4.2)
\geometry{
  a4paper,
  left=4cm,
  right=2.5cm,
  top=3cm,
  bottom=3cm,
  headheight=15pt,
  headsep=0.8cm,
  footskip=1.5cm
}

% Posisi nomor halaman sesuai diagram bidang pengetikan:
% - Header (kanan-atas) pada 1.5 cm dari tepi atas
% - Footer (halaman bab, tengah-bawah) pada 1.5 cm dari tepi bawah
% Dalam LaTeX, headsep adalah jarak header ke bidang teks, sedangkan footskip
% adalah jarak bidang teks ke footer. Keduanya diatur melalui geometry agar
% top/bottom tetap menjadi batas bidang pengetikan 3 cm sesuai panduan.

% ============================================================================
%                         TYPOGRAPHY & SPACING
% ============================================================================
% Line spacing: panduan Udinus BAB IV 4.2(k) mengatur seluruh isi tesis
% menggunakan spasi 1.5. Pengecualian berlaku untuk judul/keterangan gambar,
% tabel, grafik, dan lampiran yang boleh dibuat lebih rapat sesuai kebutuhan.
\onehalfspacing % Mengatur spasi isi utama menjadi 1.5 sesuai panduan tesis Udinus.

% Paragraph layout: paragraf lanjutan diberi indentasi, sementara paragraf
% pertama setelah heading dibiarkan mulai dari batas kiri bidang pengetikan.
% Tidak ada parskip tambahan agar jarak antarparagraf tetap mengikuti spasi 1.5,
% bukan menjadi renggang seperti gaya dokumen modern.
\setlength{\parindent}{24pt} % Indentasi paragraf lanjutan; kira-kira setara 4 karakter pada ukuran 12pt.
\setlength{\parskip}{0pt} % Tidak menambah jarak ekstra antarparagraf; jarak cukup mengikuti spasi 1.5.


% ============================================================================
%                         CHAPTER & SECTION HEADINGS
% ============================================================================
% Chapter/section formatting and spacing rules

% Chapter title: uppercase, centered at top of text area; roman numeral label
% Judul bab ala Udinus: baris 1 = "BAB ROMAWI", baris 2 = judul (uppercase)
\newcommand{\udinusChapterPrefix}{BAB}
\newcommand{\udinusChapterLabel}{\Roman{chapter}}
\titleformat{\chapter}[display]
  {\bfseries\filcenter}
  {\MakeUppercase{\udinusChapterPrefix\ \udinusChapterLabel}}
  {0pt}
  {\bfseries\MakeUppercase}

% Chapter spacing: judul bab ditempatkan dekat batas atas bidang pengetikan
% seperti contoh panduan. Jarak setelah judul bab dibuat setara 1.5 baris;
% panduan Word menyebut 3 baris, tetapi nilai literal tersebut terlalu renggang
% dalam LaTeX karena tinggi baris dan mekanisme heading berbeda.
% Format: \titlespacing*{command}{left}{before-sep}{after-sep}
\titlespacing*{\chapter}{0pt}{-1.5cm}{1.5\baselineskip} % {kiri}{sebelum}{sesudah}: judul bab dinaikkan mendekati margin atas, lalu diberi jarak 1.5 baris ke teks.

% Section and subsection: left aligned, Title Case by user; spacing adjusted for thesis readability
% Format: \titlespacing*{command}{left}{before-sep}{after-sep}
% Panduan Udinus berbasis Word menyebut 2 baris antara teks dan subjudul,
% serta antara subjudul dan teks berikutnya. Di LaTeX, 2\baselineskip literal
% terlihat terlalu renggang dan boros halaman, sehingga digunakan adaptasi yang
% tetap akademik: 1.5 baris sebelum heading dan 0.75 baris setelah heading.
\titlespacing*{\section}{0pt}{0.75\baselineskip}{0.45\baselineskip} % Section utama: 1.5 baris dari teks sebelumnya, 0.75 baris ke teks berikutnya.
\titlespacing*{\subsection}{0pt}{0.75\baselineskip}{0.45\baselineskip} % Subsection: spacing disamakan agar hierarki konsisten dan tidak terlihat menempel.
\titlespacing*{\subsubsection}{0pt}{0.75\baselineskip}{0.45\baselineskip} % Subsubsection: cukup lega untuk judul kecil seperti kategori dataset, tetapi tetap hemat halaman.

% Set heading sizes: BAB (chapter) and sub-bab use 12pt bold (no oversized sections)
\titleformat{\section}{\bfseries\normalsize}{\thesection}{0.8em}{} % Section dicetak tebal ukuran normal 12pt; jarak nomor ke judul 0.8em.
\titleformat{\subsection}{\bfseries\normalsize}{\thesubsection}{0.8em}{} % Subsection mengikuti ukuran section agar tidak terlalu besar untuk format tesis.
% Subsubsection: format huruf (a), (b), (c) — lazim PTS Indonesia level ke-4
\setcounter{secnumdepth}{3} % Menampilkan penomoran sampai level subsubsection.
\renewcommand{\thesubsubsection}{\alph{subsubsection}} % Mengubah nomor subsubsection menjadi huruf: a, b, c, dst.
\titleformat{\subsubsection}{\bfseries\normalsize}{(\thesubsubsection)}{0.8em}{} % Subsubsection tampil sebagai (a), (b), (c) dengan teks tebal ukuran normal.

% ============================================================================
%                         HEADER & FOOTER
% ============================================================================
% Headers/footers: page number at top-right; chapter-opening pages show number at bottom-center
\pagestyle{fancy}
\fancyhf{}
% Header kanan default: hanya nomor halaman
\fancyhead[R]{\thepage}
\renewcommand{\headrulewidth}{0pt}

% Redefine plain page style (used by chapter openings)
\fancypagestyle{plain}{
  \fancyhf{}
  \fancyfoot[C]{\thepage}
  \renewcommand{\headrulewidth}{0pt}
  \renewcommand{\footrulewidth}{0pt}
}

% ============================================================================
%                         DAFTAR ISI/TABEL/GAMBAR (TOC/LOT/LOF)
% ============================================================================
% Khusus untuk halaman Daftar Isi: judul dibuat rata tengah sesuai permintaan,
% sementara label "Halaman" tetap rata kanan seperti sebelumnya.
% Trik rata tengah judul menggunakan \hfill di sebelum dan sesudah judul
% (yang kedua disisipkan melalui \cftaftertoctitle di bawah).
% Gunakan font judul standar tanpa pemaksaan perataan di sini;
% perataan akan diatur lewat \cftmaketoctitle yang kita custom.
\renewcommand{\cfttoctitlefont}{\hfill\bfseries\normalsize}
\renewcommand{\cftloftitlefont}{\hfill\bfseries\normalsize}
\renewcommand{\cftlottitlefont}{\hfill\bfseries\normalsize}
% Tambahkan \hfill sebelum \par untuk menutup pasangan \hfill di depan judul
% sehingga judul "DAFTAR ISI" menjadi benar-benar di tengah.
% Baris "Halaman" tetap rata kanan.
% Seimbangkan \hfill untuk judul DAFTAR GAMBAR dan DAFTAR TABEL
\renewcommand{\cftafterloftitle}{\hfill\mbox{}\par\noindent\hfill\normalfont\normalsize Halaman\par\vspace{-0.2\baselineskip}}
\renewcommand{\cftafterlottitle}{\hfill\mbox{}\par\noindent\hfill\normalfont\normalsize Halaman\par\vspace{-0.2\baselineskip}}
% Jarak vertikal sebelum judul daftar disamakan dengan \chapter* agar posisi y
% judul DAFTAR ISI, DAFTAR GAMBAR, dan DAFTAR TABEL sejajar dengan BAB dan
% DAFTAR ISTILAH. Nilai -1.5 cm mengikuti \titlespacing*{\chapter} di atas.
\renewcommand{\cftbeforetoctitleskip}{-1.5cm}
\renewcommand{\cftaftertoctitleskip}{1\baselineskip}
\renewcommand{\cftbeforeloftitleskip}{-1.5cm}
\renewcommand{\cftafterloftitleskip}{1\baselineskip}
\renewcommand{\cftbeforelottitleskip}{-1.5cm}
\renewcommand{\cftafterlottitleskip}{1\baselineskip}
% Padatkan jeda antar entri pada daftar isi/daftar gambar/daftar tabel
\setlength{\cftbeforechapskip}{0pt}
\setlength{\cftbeforesecskip}{0pt}
\setlength{\cftbeforefigskip}{0pt}
\setlength{\cftbeforetabskip}{0pt}
% Pusatkan judul DAFTAR ISI tanpa mengubah mekanisme internal tocloft.
% Trik: \cfttoctitlefont diawali \hfill dan ditutup \hfill\mbox{} di awal
% \cftaftertoctitle agar judul benar-benar di tengah. Setelah itu, cetak
% label "Halaman" tetap rata kanan.
\renewcommand{\cftaftertoctitle}{\hfill\mbox{}\par\noindent\hfill\normalfont\normalsize Halaman\par\vspace{-0.2\baselineskip}}
\let\udinusOrigTableofcontents\tableofcontents
\renewcommand{\tableofcontents}{%
  {\begingroup
   \setstretch{1.25}%
   \udinusOrigTableofcontents
  \endgroup}}
\let\udinusOrigListoffigures\listoffigures
\renewcommand{\listoffigures}{%
  {\begingroup
   \setstretch{1.25}%
   \udinusOrigListoffigures
  \endgroup}}
\let\udinusOrigListoftables\listoftables
\renewcommand{\listoftables}{%
  {\begingroup
   \setstretch{1.25}%
   \udinusOrigListoftables
  \endgroup}}
\makeatletter
\renewcommand{\listofalgorithms}{%
  {\begingroup
   \setstretch{1.25}%
   \ifx\algocf@seclistalgo\chapter
     \if@twocolumn\@restonecoltrue\onecolumn\else\@restonecolfalse\fi
   \fi
   \algocf@seclistalgo*{\listalgorithmcfname}%
   \@mkboth{\MakeUppercase\listalgorithmcfname}{\MakeUppercase\listalgorithmcfname}%
   \par\noindent\hfill\normalfont\normalsize Halaman\par\vspace{-0.2\baselineskip}%
   \@starttoc{loa}%
   \ifx\algocf@seclistalgo\chapter
     \if@restonecol\twocolumn\fi
   \fi
  \endgroup}}


% Daftar Isi dan judul bab: pengaturan prefiks "BAB" atau "LAMPIRAN"
\renewcommand{\udinusChapterPrefix}{BAB}% Teks prefiks pada judul bab
\renewcommand{\udinusChapterLabel}{\Roman{chapter}}% Format nomor bab di judul
\newcommand{\@chapapp@toc}{BAB}% Prefiks di daftar isi
\newcommand{\udinus@setchaptertocprefix}[1]{\gdef\@chapapp@toc{#1}}
\let\udinus@orig@appendix\appendix
\renewcommand{\appendix}{%
  \udinus@orig@appendix
  \gdef\udinusChapterPrefix{LAMPIRAN}%
  \gdef\udinusChapterLabel{\Alph{chapter}}%
  \udinus@setchaptertocprefix{LAMPIRAN}%
  \addtocontents{toc}{\protect\udinus@setchaptertocprefix{LAMPIRAN}}%
  \addtocontents{toc}{\protect\setlength{\cftchapnumwidth}{7.5em}}% Perlebar kolom nomor untuk LAMPIRAN
}
\renewcommand{\cftchappresnum}{\@chapapp@toc\ }

% Command wrapper tanpa @ untuk digunakan di dokumen utama
\newcommand{\setChapterTocPrefix}[1]{\udinus@setchaptertocprefix{#1}}
\makeatother
\renewcommand{\cftchapaftersnum}{\enspace}% spasi setelah nomor bab (lebih rapat)
% Lebar kolom nomor bab di TOC agar judul lebih dekat tanpa menabrak nomor
\setlength{\cftchapnumwidth}{3.8em}
% Pastikan entri bab tidak diindent berbeda dari entri bagian
\setlength{\cftchapindent}{0pt}

% ============================================================================
%                         CAPTIONS & FLOATS
% ============================================================================
% Captions: Table caption above (default), Figure caption below (default); adjust label names done in commands.tex
\captionsetup{font=normalsize}

% Global float spacing guide: tweak these lengths to control gap between figures/tables and surrounding text.
% Float placement specifiers for e.g. \begin{figure}[tbp]:
% t = place at top of page
% b = place at bottom of page
% h = try placing near the "here" location in source
% p = allow placement on a float-only page
\setlength{\floatsep}{12pt plus 2pt minus 2pt}
\setlength{\textfloatsep}{12pt plus 2pt minus 2pt}
\setlength{\intextsep}{12pt plus 2pt minus 2pt}

% Force floats on float-only pages to start at the top (not centered vertically).
% Default LaTeX centers floats on \floatpage; these settings push them to the top.
\makeatletter
\setlength{\@fptop}{0pt}           % No extra glue at top of float page
\setlength{\@fpsep}{8pt plus 1fil} % Space between floats on a float page
\setlength{\@fpbot}{0pt plus 1fil} % Minimal glue at bottom
\makeatother

% Relax float placement limits so [t] figures are more likely to be placed at top.
\renewcommand{\topfraction}{0.9}      % Up to 90% of a text page can be floats at top
\renewcommand{\textfraction}{0.1}     % At least 10% of page must be text
\renewcommand{\floatpagefraction}{0.8} % A float page must be at least 80% full

% ============================================================================
%                         PSEUDOCODE/ALGORITMA (Algorithm2e)
% ============================================================================
\SetKwInput{KwInput}{Input}
\SetKwInput{KwOutput}{Output}
\SetKw{KwAnd}{and}
\SetKw{KwOr}{or}
\SetAlgoCaptionSeparator{.}
\SetAlgoNlRelativeSize{-1}
\SetAlFnt{\small}
\SetKwComment{Comment}{$\triangleright$ }{}
\SetAlgoSkip{bigskip} % Tambahan jarak vertikal sebelum/sesudah algoritma

% Penomoran algoritme per bab: x.y (IEEE style)
\renewcommand{\thealgocf}{\arabic{chapter}.\arabic{algocf}}
\makeatletter
\@addtoreset{algocf}{chapter}
\makeatother

% Nama float - Gaya Indonesia
\SetAlgorithmName{Algoritma}{Algoritma}{DAFTAR ALGORITMA}
